\documentclass[11pt]{article}
\usepackage{amsmath, amsthm, amssymb}
\usepackage[margin=1in]{geometry}
\usepackage{hyperref}
\usepackage{array}

\newtheorem{theorem}{Theorem}
\newtheorem{proposition}[theorem]{Proposition}
\newtheorem{lemma}[theorem]{Lemma}
\newtheorem{corollary}[theorem]{Corollary}
\theoremstyle{definition}
\newtheorem{definition}[theorem]{Definition}
\newtheorem{conjecture}[theorem]{Conjecture}
\newtheorem{observation}[theorem]{Observation}
\theoremstyle{remark}
\newtheorem*{remark}{Remark}

\newcommand{\bbZ}{\mathbb{Z}}
\newcommand{\bbQ}{\mathbb{Q}}
\newcommand{\bbN}{\mathbb{N}}
\newcommand{\Pq}{\Pi_q}
\newcommand{\Hq}{\mathcal{H}_q(S_n)}
\newcommand{\Tw}{T_w}

\title{A new $(a,b)$-class atom $h_{(7,2)}(q)$\\
       and the first multi-shape $\sigma_1$-coincidence at $n=8$}
\author{Clio\\\small (with thanks to Robin Langer)}
\date{2026-04-30}

\begin{document}
\maketitle

\begin{abstract}
For each partition $\lambda$ of $n$ let $\sigma_1(V_\lambda)(q)$ denote the
trace of the staircase Hecke product
$\Pq^{(n)} = (1+T_1)(1+T_2)(1+T_1)(1+T_3)(1+T_2)(1+T_1)\cdots(1+T_{n-1})\cdots(1+T_1)$
acting on the Specht module $V_\lambda$ via Young's seminormal form.
Empirically
\[
\sigma_1(V_\lambda)(q) \;=\; q^{a_\lambda}(1+q)^{b_\lambda}\,P_\lambda(q),
\]
with $P_\lambda$ palindromic and $(a_\lambda, b_\lambda)$ maximal among
factorizations of this shape.  Define the \emph{$(a,b)$-class} of $\lambda$
as the equivalence class of partitions sharing $(a_\lambda, b_\lambda)$.

We exhibit the first multi-element $(a,b)$-class:
$\lambda = (3,3,2)$ and $\mu = (4,2,1,1)$ are both partitions of $8$, both
self-conjugate, both have $(a,b) = (7,2)$, and they have the
\emph{identical} $\sigma_1$ polynomial despite
$\dim V_\lambda = 42 \neq 90 = \dim V_\mu$.
The common factor $h_{(7,2)}(q)$ is irreducible palindromic of degree
$12$ with leading coefficient $1$; it is novel --- not in any previously
catalogued atom-alphabet position.  We rule out W-graph local structure as
an explanation by direct comparison of $\tau$-multisets, and reframe the
coincidence as a Markov-trace cancellation
$\sum_\nu |C_\nu|\,P_\nu(q)\,\delta(c_\nu) \equiv 0$ in $\bbQ[q]$, with
$\delta = \chi_\lambda - \chi_\mu$ supported on $8$ of $22$ conjugacy
classes including the identity, where $\delta$ has defect
$-48 = \dim V_\lambda - \dim V_\mu$.
\end{abstract}

\section{Introduction}

Let $\Hq$ denote the Iwahori--Hecke algebra of type $A_{n-1}$ over
$\bbQ(q)$, with generators $T_1, \dots, T_{n-1}$ and quadratic relation
$T_i^2 = (q-1)T_i + q$.  Fix the staircase reduced word
\[
\mathbf{w}_0
= (s_1)(s_2 s_1)(s_3 s_2 s_1) \cdots (s_{n-1}\cdots s_1)
\]
for the longest element $w_0 \in S_n$.  Define the staircase Hecke
product
\[
\Pq^{(n)}
\;=\; \prod_{i=1}^{n-1}\prod_{j=i}^{1}(1+T_j)
\;=\; (1+T_1)(1+T_2)(1+T_1)(1+T_3)(1+T_2)(1+T_1)\cdots
\]
on $\binom{n}{2}$ factors.  For each partition $\lambda \vdash n$ let
$V_\lambda$ be the Specht module realized via Young's seminormal form,
and define the $\sigma_1$-trace
\[
\sigma_1(V_\lambda)(q) \;=\; \mathrm{tr}\bigl(\Pq^{(n)} \,\big|\, V_\lambda\bigr)
\;\in\; \bbZ[q].
\]

\begin{definition}[$(a,b)$-class]
The unique factorization
\(
\sigma_1(V_\lambda)(q) = q^{a_\lambda}(1+q)^{b_\lambda} P_\lambda(q)
\)
with $P_\lambda(0) \neq 0$, $P_\lambda(-1) \neq 0$, and $P_\lambda$
palindromic, defines the pair $(a_\lambda, b_\lambda) \in \bbN^2$.
We call this pair the \emph{$(a,b)$-degree} of $\lambda$, and the
equivalence relation $\lambda \sim \mu \iff (a_\lambda, b_\lambda) =
(a_\mu, b_\mu)$ partitions $\{\lambda \vdash n\}$ into
\emph{$(a,b)$-classes}.
\end{definition}

We exclude rank-isolated shapes ($\ell(\lambda) > (n+1)/2$, where
$\Pq$ acts as zero on $V_\lambda$) from the enumeration.

\begin{observation}[Main observation]
At $n=7$, every $(a,b)$-class is a singleton: there are $11$
non-zero classes, all of size $1$.  At $n=8$, there are $14$
non-zero $(a,b)$-classes, exactly one of which has size $2$:
\[
\bigl\{(3,3,2),\,(4,2,1,1)\bigr\}, \qquad (a,b) = (7,2).
\]
Moreover $\sigma_1(V_{(3,3,2)})(q) = \sigma_1(V_{(4,2,1,1)})(q)$ as
polynomials in $q$, despite $\dim V_{(3,3,2)} = 42$ and
$\dim V_{(4,2,1,1)} = 90$.  The common polynomial $P(q)$ is irreducible
palindromic of degree $12$ over $\bbQ$.
\end{observation}

The remainder of this note documents and partially structures this
finding.

\section{The Rosetta Stone factorization}

The factorization
$\sigma_1(V_\lambda) = q^{a_\lambda}(1+q)^{b_\lambda} P_\lambda(q)$ is
the per-irrep instance of what we have been calling the
\emph{Rosetta Stone}: the trace of $\Pq^{(n)}$ on each Specht module
splits cleanly into
\begin{itemize}
\item a $q$-power $q^{a_\lambda}$ (Lusztig $a$-value, content shift),
\item a $(1+q)^{b_\lambda}$ factor (Soergel / categorical
      multiplicity), and
\item a palindromic remainder $P_\lambda(q)$ (the \emph{atom}).
\end{itemize}
This per-irrep factorization was first verified explicitly at $n=4$
on April~22, 2026 (cf.\ memory entry
\texttt{rosetta\_stone\_n4\_verified}); the present work confirms it
through $n=8$ and isolates the new degree-12 atom $h_{(7,2)}(q)$
described below.

For all $\lambda \vdash 8$ that are not rank-isolated, $P_\lambda(q)$
is palindromic in our computations; this is consistent with the
W-graph parity-sheaf self-duality predicted by the monodromic Hecke
hypothesis.

\section{$(a,b)$-class enumeration at $n=7$ and $n=8$}

We computed $\sigma_1(V_\lambda)(q)$ for every non-rank-isolated
$\lambda \vdash 7, 8$ via Lagrange interpolation from $32$ exact
$\bbQ$-evaluations of the seminormal-form action of $\Pq^{(n)}$ on
$V_\lambda$.  Table~\ref{tab:n7} reports the $(a,b)$-degrees for
$n=7$; Table~\ref{tab:n8} for $n=8$.  Every $P_\lambda$ in both
tables is palindromic.

\begin{table}[ht]
\centering
\begin{tabular}{lrrrr}
\hline
$\lambda$ & $\dim V_\lambda$ & $a$ & $b$ & $\deg P$ \\
\hline
$(2,2,2,1)$  & 14 & 9 & 3 &  0 \\
$(3,2,1,1)$  & 35 & 7 & 1 &  6 \\
$(3,2,2)$    & 21 & 6 & 1 &  8 \\
$(3,3,1)$    & 21 & 5 & 1 & 10 \\
$(4,1,1,1)$  & 20 & 6 & 3 &  6 \\
$(4,2,1)$    & 35 & 4 & 1 & 12 \\
$(4,3)$      & 14 & 3 & 3 & 12 \\
$(5,1,1)$    & 15 & 3 & 5 & 10 \\
$(5,2)$      & 14 & 2 & 5 & 12 \\
$(6,1)$      &  6 & 1 &11 &  8 \\
$(7)$        &  1 & 0 &21 &  0 \\
\hline
\end{tabular}
\caption{$(a,b)$-degrees at $n=7$, non-rank-isolated shapes.  All
$11$ classes are singletons.}
\label{tab:n7}
\end{table}

\begin{table}[ht]
\centering
\begin{tabular}{lrrrr}
\hline
$\lambda$ & $\dim V_\lambda$ & $a$ & $b$ & $\deg P$ \\
\hline
$(2,2,2,2)$    & 14 & 12 & 4 &  0 \\
$(3,2,2,1)$    & 70 &  9 & 2 &  8 \\
$(3,3,1,1)$    & 56 &  8 & 2 & 10 \\
$\mathbf{(3,3,2)}$    & \textbf{42} &  \textbf{7} & \textbf{2} & \textbf{12} \\
$\mathbf{(4,2,1,1)}$  & \textbf{90} &  \textbf{7} & \textbf{2} & \textbf{12} \\
$(4,2,2)$      & 56 &  6 & 2 & 14 \\
$(4,3,1)$      & 70 &  5 & 2 & 16 \\
$(4,4)$        & 14 &  4 & 4 & 16 \\
$(5,1,1,1)$    & 35 &  6 & 4 & 12 \\
$(5,2,1)$      & 64 &  4 & 2 & 18 \\
$(5,3)$        & 28 &  3 & 4 & 18 \\
$(6,1,1)$      & 21 &  3 & 8 & 14 \\
$(6,2)$        & 20 &  2 & 8 & 16 \\
$(7,1)$        &  7 &  1 &16 & 10 \\
$(8)$          &  1 &  0 &28 &  0 \\
\hline
\end{tabular}
\caption{$(a,b)$-degrees at $n=8$, non-rank-isolated shapes.  The
unique multi-element $(a,b)$-class
$\{(3,3,2), (4,2,1,1)\}$ is bolded.}
\label{tab:n8}
\end{table}

\section{The polynomial $h_{(7,2)}(q)$}

\begin{theorem}[Common atom]
$\sigma_1(V_{(3,3,2)})(q) = \sigma_1(V_{(4,2,1,1)})(q) = q^{7}(1+q)^{2}\,
h_{(7,2)}(q)$, where
\begin{align*}
h_{(7,2)}(q) \;=\;& 1 + 27q + 325q^{2} + 2169q^{3} + 8428q^{4} + 19143q^{5} \\
&{}+ 25218q^{6} \\
&{}+ 19143q^{7} + 8428q^{8} + 2169q^{9} + 325q^{10} + 27q^{11} + q^{12}.
\end{align*}
The polynomial $h_{(7,2)}$ is palindromic of degree $12$, has leading
and constant coefficient $1$, and is irreducible over $\bbQ$.
\end{theorem}

\begin{proof}[Computational verification]
Both $\sigma_1$-polynomials are computed by exact Lagrange
interpolation in $\bbQ[q]$ from $32$ evaluations of
$\mathrm{tr}(\Pq^{(8)} | V_\lambda)$ via Young's seminormal form
(matrix size $42 \times 42$ for $\lambda = (3,3,2)$ and $90 \times 90$
for $\lambda = (4,2,1,1)$); both interpolations return the same
$22$-coefficient polynomial $q^7 (1+q)^2 h(q)$ with the coefficients
displayed above.

Palindromy: the coefficient list
$(1, 27, 325, 2169, 8428, 19143, 25218, 19143, 8428, 2169, 325, 27, 1)$
is symmetric, so $q^{12} h_{(7,2)}(1/q) = h_{(7,2)}(q)$.

Irreducibility: a SymPy factorization confirms that $h_{(7,2)}(q)$ is
irreducible over $\bbQ$.  No cyclotomic polynomial divides it (a quick
check evaluates $h_{(7,2)}$ at all primitive roots of unity of order
$\le 13$).  No previously catalogued atom of degree $\le 12$ from the
$n \le 7$ alphabet (see the memory entry \texttt{atom\_alphabet}) is a
factor.
\end{proof}

The values
\[
h_{(7,2)}(0) = 1, \qquad
h_{(7,2)}(1) = 85404 = 2^{2} \cdot 3 \cdot 11 \cdot 647,
\]
are noteworthy: the prime factor $647$ appears in no previously
catalogued atom value $h(1)$, and the central coefficient $25218 =
2 \cdot 3^{2} \cdot 1401$ does not match any Eulerian, Narayana, or
derangement-RTL value at the relevant length.  We therefore record
$h_{(7,2)}(q)$ as a \emph{new entry} in the atom alphabet, distinct
from every entry isolated through $n=7$.

The $\sigma_1$ value at $q=1$ is
\[
\sigma_1(V_{(3,3,2)})(1) = \sigma_1(V_{(4,2,1,1)})(1) = 1^{7} \cdot 2^{2}
\cdot 85404 = 341616 = 2^{4} \cdot 3 \cdot 11 \cdot 647.
\]

\section{Self-conjugacy and Lusztig $a$-value}

Both $\lambda = (3,3,2)$ and $\mu = (4,2,1,1)$ are self-conjugate
($\lambda = \lambda'$, $\mu = \mu'$).  Recall that $n(\lambda) =
\sum_{i \ge 1} (i-1)\lambda_i$ is Lusztig's $a$-invariant for type $A$.
Direct computation gives
\[
n(\lambda) = 0\cdot 3 + 1\cdot 3 + 2\cdot 2 = 7, \qquad
n(\mu)     = 0\cdot 4 + 1\cdot 2 + 2\cdot 1 + 3\cdot 1 = 7,
\]
and by self-conjugacy $n(\lambda) = n(\lambda')$, $n(\mu) = n(\mu')$,
so $\lambda$ and $\mu$ have the \emph{same Lusztig $a$-value}, namely
$a(\lambda) = a(\mu) = 7$.  This matches our
$a_{\lambda} = a_{\mu} = 7$ in the $(a,b)$-degree.  In particular the
$q$-power exponent in the Rosetta Stone factorization equals the
Lusztig $a$-value here.  Both shapes therefore have the same Lusztig
two-sided cell-degree contribution; this is a necessary condition for
the $\sigma_1$-coincidence but, as the next section shows, far from
sufficient.

\section{W-graph local structure does not explain the coincidence}

For each standard Young tableau $T$ of shape $\lambda \vdash n$ let
$\mathrm{Des}(T) = \{i \in \{1, \dots, n-1\} : i+1 \text{ lies strictly below } i \text{ in } T\}$
be its descent set.  The multiset of descent sets $\{\mathrm{Des}(T)
: T \in \mathrm{SYT}(\lambda)\}$ records the local W-graph structure
attached to the Kazhdan--Lusztig left cell of shape $\lambda$.

\begin{proposition}
The descent multisets of $(3,3,2)$ and $(4,2,1,1)$ differ strongly.
Even after dimension-normalization:
\begin{itemize}
\item $\lambda = (3,3,2)$: descent-set sizes are
  $\{2, 3, 4, 5\}$ with normalized weights
  $(\tfrac{1}{14}, \tfrac{3}{7}, \tfrac{3}{7}, \tfrac{1}{14})$;
\item $\mu = (4,2,1,1)$: descent-set sizes are $\{3, 4\}$ with
  normalized weights $(\tfrac{1}{2}, \tfrac{1}{2})$.
\end{itemize}
In particular $\lambda$ has $3$ tableaux with descent set of size $5$;
$\mu$ has none. $\mu$ has $33$ tableaux with descent set of size $3$
that do not appear among the descent sets of $\lambda$.  The
intersection of the descent-set supports has $33$ elements out of
$|\lambda|=35$ and $|\mu|=75$.
\end{proposition}

\begin{proof}
Direct enumeration over $\mathrm{SYT}((3,3,2))$ (42 tableaux) and
$\mathrm{SYT}((4,2,1,1))$ (90 tableaux); see the script
\texttt{tau\_multisets.py} in the supporting computation directory.
\end{proof}

\begin{corollary}
The $\sigma_1$-coincidence is not visible at the W-graph descent
level.  In particular it is not a shadow of any local
descent-multiset agreement, and therefore is not predicted by any
parity-sheaf or local KL-graph isomorphism between the two cells.
\end{corollary}

\section{Markov-trace reframing}

Decompose $\Pq^{(n)} \in \Hq$ in the $T_w$ basis,
$\Pq^{(n)} = \sum_{w \in S_n} c_w(q)\,T_w$, and let
$\widetilde{P}_\nu(q) = \sum_{w \in C_\nu} c_w(q)$
where $C_\nu \subset S_n$ is the conjugacy class of cycle type $\nu$.
Each Hecke character $\chi_\lambda^{q}$ evaluates linearly on $T_w$;
projecting to ordinary characters at $q \to 1$ yields
\[
\sigma_1(V_\lambda)(q) \;=\; \sum_{\nu \vdash n} P_\nu(q)\,\chi_\lambda(c_\nu),
\]
for an appropriate set of class polynomials $P_\nu(q)$ adapted to the
Geck--Pfeiffer character formula on the Hecke algebra.  In particular
\[
\sigma_1(V_\lambda)(q) - \sigma_1(V_\mu)(q)
\;=\; \sum_{\nu \vdash n} P_\nu(q)\,\delta(c_\nu),
\qquad \delta = \chi_\lambda - \chi_\mu.
\]

\begin{proposition}[Support of $\delta$ at $n = 8$]
For $\lambda = (3,3,2)$ and $\mu = (4,2,1,1)$, the difference
$\delta = \chi_\lambda - \chi_\mu$ is supported on exactly $8$ of
the $22$ conjugacy classes of $S_8$:
\[
\begin{array}{lrrr}
\nu & \chi_\lambda(c_\nu) & \chi_\mu(c_\nu) & \delta(c_\nu) \\
\hline
(7,1)             &  0 & -1 &  1 \\
(5,3)             & -1 &  0 & -1 \\
(5,1,1,1)         &  2 &  0 &  2 \\
(4,2,1,1)         & -2 &  2 & -4 \\
(3,2,2,1)         &  2 &  0 &  2 \\
(3,1,1,1,1,1)     & -6 &  0 & -6 \\
(2,2,1,1,1,1)     &  2 & -6 &  8 \\
(1^{8})           & 42 & 90 & -48
\end{array}
\]
On the other $14$ conjugacy classes both characters agree.
\end{proposition}

\begin{proof}
Murnaghan--Nakayama recursion via $\beta$-numbers, performed in
exact integer arithmetic, with the sanity check $\sum_\lambda
\dim(V_\lambda)^{2} = 8! = 40320$ passing; full character table on
the two rows is recorded in
\texttt{2026-04-30-tcore-test/RESULT.md}.
\end{proof}

\begin{corollary}
The $\sigma_1$-coincidence
$\sigma_1(V_{(3,3,2)})(q) = \sigma_1(V_{(4,2,1,1)})(q)$ in $\bbZ[q]$
is equivalent to the Markov-trace identity
\[
\sum_{\nu \in \mathrm{supp}(\delta)} P_\nu(q)\,\delta(c_\nu) \;\equiv\; 0
\quad \text{in } \bbQ[q].
\]
At $q = 1$, this becomes a non-trivial linear identity among the eight
values $\{1, -1, 2, -4, 2, -6, 8, -48\}$ weighted by the $q=1$ values
of the eight class polynomials.
\end{corollary}

\begin{remark}
The defect at the identity, $\delta(c_{(1^{8})}) = -48 = 42 - 90$,
must be exactly cancelled by contributions from the other seven
classes in $\mathrm{supp}(\delta)$.  Three of those classes ---
$(4,2,1,1)$, $(2,2,1,1,1,1)$, and $(1^8)$ --- have both
$\chi_\lambda$ and $\chi_\mu$ nonzero, so the identity is genuinely
$q$-weighted: it cannot be reduced to an obvious character-zero
mechanism (such as a $t$-core / Type-II coincidence).
\end{remark}

\section{Conjecture: every multi-element $(a,b)$-class carries an atom}

\begin{conjecture}[Class atom]
For every $n$ and every multi-element $(a,b)$-class
$[\lambda] \subset \{\nu \vdash n : (a_\nu, b_\nu) = (a,b)\}$, the
greatest common divisor
\[
h_{[\lambda]}(q) \;:=\; \gcd_{\mu \in [\lambda]} P_\mu(q)
\]
is a non-trivial palindromic polynomial in $\bbZ[q]$, the
\emph{$(a,b)$-class atom}.  Moreover $h_{[\lambda]}$ enters the atom
alphabet as a new entry that is not shared with any singleton class
$P_\nu$ at the same $n$.
\end{conjecture}

The polynomial $h_{(7,2)}(q)$ above is the first instance of a class
atom: the $(a,b)$-class $\{(3,3,2), (4,2,1,1)\}$ at $n = 8$ has
$P_{(3,3,2)} = P_{(4,2,1,1)} = h_{(7,2)}$, so the gcd is the whole
common $P$.  At $n \le 7$ no multi-element class exists, so the
conjecture has no content there.

\section{Open questions}

\begin{enumerate}
\item \textbf{Hecke class polynomials.}  Compute the Geck--Pfeiffer
class polynomials $P_\nu(q)$ for $\Pq^{(8)}$ explicitly, and verify
the cancellation identity $\sum_\nu |C_\nu|\,P_\nu(q)\,\delta(c_\nu)
\equiv 0$ in $\bbQ[q]$ directly.  This requires expanding
$\Pq^{(8)}$ in the $T_w$ basis ($28$-letter staircase word, of order
$10^{6}$ terms in $\bbZ[q]$); feasible in plain Python.
\item \textbf{Categorical home of $h_{(7,2)}$.}  The atom alphabet
isolated through $n=7$ has explained entries (Eulerian $E_2$,
Narayana, derangement-RTL, $A_9, A_{10}$, the orphan $A_{11}$ from
\texttt{novel\_atom\_deg10}).  Where does $h_{(7,2)}$ live?  One
candidate is a singular-Soergel atomic Leibniz contribution in the
Patimo categorical framework: the $(a,b) = (7,2)$ cell $S_8$ admits
two cells of the same Lusztig $a$-value but distinct
two-sided structure, which could produce a single Hochschild--type
contribution shared between them.
\item \textbf{Higher $n$.}  At $n = 9, 10$, how many multi-element
$(a,b)$-classes appear, and do they all have non-trivial class
atoms?  Is the multiplicity of multi-element classes monotone in
$n$?  Are class-atom polynomials closed under $n \mapsto n+1$
branching, or do they form an independent layer of atoms?
\item \textbf{Self-conjugacy.}  Both members of the $(7,2)$ class at
$n=8$ are self-conjugate.  Is this generic, i.e.\ do all
multi-element $(a,b)$-classes consist entirely of self-conjugate
shapes?  If so, this restricts the conjectural class-atom alphabet
to a controlled subset of irreducible palindromes.
\end{enumerate}

\paragraph{Acknowledgements.} Thanks to Robin Langer for sustained
attention to the staircase Hecke trace program and for steady
calibration of speculative versus solid claims throughout the past
month.

\paragraph{Computational data.} Source code and exact polynomial
values are recorded under
\verb|/home/clio/projects/scratch/2026-04-30-ab-classes-n78/| and
\verb|/home/clio/projects/scratch/2026-04-30-sigma1-equality-structure/|;
the Murnaghan--Nakayama character verification of the
$\delta$-support is recorded under
\verb|/home/clio/projects/scratch/2026-04-30-tcore-test/|.

\end{document}
